\section*{Hoofdstuk \ref{ch:g12:plant-rooted-life} --- Het leven van de plant: gedijen met wortels}

\begin{solution}{exo:g12:plant-rooted-life:1}
Het wortelstelsel, waarvan de \omterm{prop:g12:plant-rooted-life:surfaces}{wortelharen} een uitgestrekt oppervlak aan de
bodem bieden, en het spruitstelsel, waarvan de bladeren een uitgestrekt
oppervlak aan licht en lucht bieden.
\end{solution}

\begin{solution}{exo:g12:plant-rooted-life:2}
Een porie in de huid van het blad, begrensd door twee sluitcellen.
Koolstofdioxide komt binnen, waterdamp (en zuurstof) gaat naar buiten.
Open in het licht, dicht 's nachts en bij droogte.
\end{solution}

\begin{solution}{exo:g12:plant-rooted-life:3}
Ruw sap: water en minerale ionen, in het \omterm{prop:g12:plant-rooted-life:saps}{xyleem}, van de wortels naar de
bladeren, getrokken door de verdamping. Bereid sap: water en suikers, in
het \omterm{prop:g12:plant-rooted-life:saps}{floëem}, van de bladeren naar de organen die suiker gebruiken of
opslaan, in beide richtingen, geduwd door de druk van de met suiker
geladen \omterm{def:g10:cells-common-unit:cell}{cellen}.
\end{solution}

\begin{solution}{exo:g12:plant-rooted-life:4}
Bouwkundig: cuticula en bast, doorns (ook brandharen, kiezelzuur,
\omterm{def:g10:cells-common-unit:organelle}{celwanden}). Scheikundig: tannines, alkaloïden (ook giffen, harsen,
melksap).
\end{solution}

\begin{solution}{exo:g12:plant-rooted-life:5}
Een gerichte groeireactie op een prikkel (licht, zwaartekracht, water).
Het auxine, gemaakt in de groeiende top van de scheut.
\end{solution}

\begin{solution}{exo:g12:plant-rooted-life:6}
$\qty{40}{cm^2} = \qty{4000}{mm^2}$: $4000 \times 200 = \num{8e5}$
\omterm{def:g12:plant-rooted-life:stomata}{huidmondjes} per blad; $\num{1.6e11}$ op de boom.
\end{solution}

\begin{solution}{exo:g12:plant-rooted-life:7}
Rond het middaguur tot de vroege namiddag. In de hitte van de namiddag
sluiten de sluitcellen de \omterm{def:g12:plant-rooted-life:stomata}{huidmondjes} deels om het waterverlies te
beperken, en de verdamping volgt de opening, niet het licht.
\end{solution}

\begin{solution}{exo:g12:plant-rooted-life:8}
De bladeren blijven groen en van water voorzien (het \omterm{prop:g12:plant-rooted-life:saps}{xyleem}, binnen in
het hout, is heel gebleven). De wortels, afgesneden van de suiker van de
bladeren, putten hun voorraden in enkele maanden uit en sterven, en de
boom met hen. Vlak boven de ring zwelt de bast op met de suiker die er
niet langer voorbij kan.
\end{solution}

\begin{solution}{exo:g12:plant-rooted-life:9}
Ongeschonden: de reactie bestaat. Top bedekt: de top is waar het licht
wordt waargenomen. Voet bedekt: de voet hoeft niet belicht te zijn om te
buigen --- hij reageert op een signaal, niet op licht. Top verwijderd: de
top is nodig; het signaal komt daarvandaan. Samen scheiden zij het
waarnemen (top) van het reageren (voet) en laten zij zien dat er iets
tussen hen reist.
\end{solution}

\begin{solution}{exo:g12:plant-rooted-life:10}
Van de zon: de verdamping van water aan het bladoppervlak, aangedreven
door zonnewarmte en droge lucht, trekt de aaneengesloten waterkolom het
\omterm{prop:g12:plant-rooted-life:saps}{xyleem} op. De plant levert de buizen en de open \omterm{def:g12:plant-rooted-life:stomata}{huidmondjes}, niet de
energie.
\end{solution}

\begin{solution}{exo:g12:plant-rooted-life:11}
Windbloem: klein, dof, geurloos, zonder nectar, enorme hoeveelheden licht
stuifmeel, vederachtige stempels --- de wind is blind en gratis, dus zet
de plant op hoeveelheid in. Insectenbloem: groot, gekleurd, geurend, met
nectar, kleverig stuifmeel in bescheiden hoeveelheden --- de drager moet
worden aangetrokken en betaald, en levert het stuifmeel nauwkeurig af.
\end{solution}

\begin{solution}{exo:g12:plant-rooted-life:12}
Zij wint water, waardoor haar \omterm{def:g10:cells-common-unit:cell}{cellen} gespannen en in leven blijven; zij
verliest de toevoer van koolstofdioxide, en daarmee de \omterm{def:g10:cell-metabolism:autotrophy}{fotosynthese}, en de
koeling door de verdamping. Gesloten \omterm{def:g12:plant-rooted-life:stomata}{huidmondjes} betekenen geen suiker:
de plant leeft van haar voorraden, en een lange droogte laat haar
verhongeren voordat zij haar uitdroogt.
\end{solution}

\begin{solution}{exo:g12:plant-rooted-life:13}
In het blad een verdediging: een gif tegen planteneters, in het weefsel
ingebouwd. In de nectar een inhuring: een lage dosis die de bestuiver laat
onthouden en terugkomen, naast de suiker betaald. Eén stof, twee
behoeften, per plaats gedoseerd.
\end{solution}

\begin{solution}{exo:g12:plant-rooted-life:14}
Bloemen met iets langere buizen worden alleen bestoven door nachtvlinders
die de bodem bereiken, dus komt hun stuifmeel bij bloemen van hun eigen
\omterm{def:g10:biodiversity-scales:species}{soort}; nachtvlinders met iets langere tongen krijgen nectar die anderen
niet krijgen; elk klein voordeel wordt geselecteerd, en buis en tong
worden over vele generaties samen langer. Verdwijnt de nachtvlinder, dan
zet de bloem geen zaad; verdwijnt de bloem, dan verliest de nachtvlinder
haar enige voedsel.
\end{solution}

\begin{solution}{exo:g12:plant-rooted-life:15}
Een plant neemt licht waar met de top van haar scheuten en de
zwaartekracht met haar wortels; zij reageert door naar het ene toe en
langs het andere te groeien; zij verdedigt zich met bast en doorns, en
maakt binnen enkele uren na een aanval giffen en verteringsblokkers; en
zij huurt insecten in om haar stuifmeel te dragen, vogels om haar zaden te
dragen en wespen om haar rupsen te doden, en betaalt in nectar, vrucht en
geur. Zij handelt zonder te bewegen.
\end{solution}

\begin{solution}{pb:g12:plant-rooted-life:1}
\textbf{1.} $\num{250000} \times \qty{25}{cm^2} = \num{6.25e6}\,
\unit{cm^2} = \qty{625}{m^2}$.

\textbf{2.} $\qty{625}{m^2} = \num{6.25e8}\,\unit{mm^2}$; $\times 250
\approx \num{1.6e11}$ \omterm{def:g12:plant-rooted-life:stomata}{huidmondjes}.

\textbf{3.} $625/150 \approx 4$ vierkante meter blad per vierkante meter
grond.

\textbf{4.} Ongeveer \qty{1000}{m^2} \omterm{prop:g12:plant-rooted-life:surfaces}{wortelharen} tegen \qty{625}{m^2}
blad: vergelijkbaar. Beide zijn uitwisselingsoppervlakken die op dezelfde
stromen zijn afgestemd --- het water dat de bladeren verliezen, moet door
de wortels worden opgenomen.

\textbf{5.} Licht wordt binnen een fractie van een millimeter weefsel
opgenomen, en koolstofdioxide verspreidt zich alleen over korte
afstanden: een dik orgaan zou de meeste van zijn \omterm{def:g10:cells-common-unit:cell}{cellen} in het donker en
zonder voedsel laten. Dunne, gestapelde bladeren brengen elke \omterm{def:g10:cells-common-unit:cell}{cel} dicht
bij licht en lucht.

\textbf{6.} $625 \times 0.6 = \qty{375}{L}$.

\textbf{7.} $625 \times 8 = \qty{5}{kg}$ koolstofdioxide, wat ongeveer
$5/1.47 \approx \qty{3.4}{kg}$ suiker geeft (of eenvoudiger: de suiker
gemaakt uit \qty{5}{kg} $\mathrm{CO_2}$ is ongeveer \qty{3.4}{kg});
gebruikt water $3.4/1.6 \approx \qty{2.1}{kg}$: ongeveer 0.6\% van de
verdampte \qty{375}{L}.

\textbf{8.} Koolstofdioxide komt binnen door dezelfde open \omterm{def:g12:plant-rooted-life:stomata}{huidmondjes}
die het water naar buiten laten; die sluiten stopt de \omterm{def:g10:cell-metabolism:autotrophy}{fotosynthese}. Het
water is de prijs van de suiker.

\textbf{9.} De zon, die water aan de bladeren verdampt. \qty{375}{kg}
\qty{20}{m} omhoogtillen is ongeveer \qty{75}{kJ} arbeid --- op zichzelf
weinig, maar een biologische pomp zou weefsels, ATP en onderhoud vergen
die de boom nu niet hoeft te bouwen.

\textbf{10.} Hij heeft zijn \omterm{def:g12:plant-rooted-life:stomata}{huidmondjes} deels gesloten tegen de droge
hitte; de toevoer van koolstofdioxide en de suikerproductie zakken met een
vergelijkbaar deel --- zo'n \qty{2}{kg} suiker verloren over de namiddag.

\textbf{11.} Ongeveer \qty{3.4}{kg} gemaakt; \qty{2}{kg} (60\%) naar de
wortels, de stam en de scheuten, \qty{0.5}{kg} (15\%) naar de eikels.

\textbf{12.} Het \omterm{prop:g12:plant-rooted-life:saps}{floëem}, van de bladeren (bron) naar de wortels, de stam,
de scheuten en de eikels (putten).

\textbf{13.} Opgeslagen als zetmeel in de stam en de wortels; het volgende
voorjaar wordt het weer in suiker omgezet om de nieuwe bladeren te bouwen
voordat die zichzelf kunnen voeden.

\textbf{14.} De suiker bereikt nog altijd de scheuten en de eikels boven
de ring; de wortels en de onderste stam krijgen niets: hun deel van
\qty{2}{kg} zakt tot wat er was opgeslagen, en zij verhongeren in de loop
van het seizoen.

\textbf{15.} De overgebleven bladeren blijven suiker maken, het
opgeslagen zetmeel van de stam en de wortels dekt het tekort, en de boom
kan vanuit zijn vele groeipunten nieuwe bladeren uitlopen; de eikels zijn
een put die de boom nog kan vullen als de suiker van de zomer volstaat.

\textbf{16.} Kleine, groenige bloemen zonder bloembladen, geur of nectar,
hangend in katjes die wolken licht stuifmeel laten vallen; grote
vederachtige stempels om het op te vangen.

\textbf{17.} Eikels die onder de ouder vallen, groeien, als zij al
groeien, in haar schaduw en beconcurreren haar; een gaai draagt ze
honderden meters ver en begraaft ze op de juiste diepte. De meeste
kwijtraken is de prijs voor er een paar goed te plaatsen.

\textbf{18.} $3000 \times 0.05 = 150$ vergeten eikels; $150/50 = 3$ jonge
bomen.

\textbf{19.} Drie jonge bomen per jaar over enkele tientallen jaren met
goede oogsten geven misschien honderd jonge bomen in het leven van de
boom, waarvan er één volwassen moet worden. De getallen zijn groot omdat
vrijwel elk zaad en elke jonge boom verloren gaat.

\textbf{20.} \qty{625}{m^2} blad en zo'n $\num{1.6e11}$ \omterm{def:g12:plant-rooted-life:stomata}{huidmondjes};
ongeveer \qty{375}{L} getild op een dag voor \qty{3.4}{kg} suiker; de wind
voor zijn stuifmeel en de gaai voor zijn eikels.
\end{solution}
